\documentclass[11pt,a4paper]{article}

% ============================================================================
% ASTRA Plotter - User Manual
% ============================================================================

\usepackage[UTF8]{ctex}
\usepackage{amsmath,amssymb}
\usepackage{graphicx}
\usepackage{hyperref}
\usepackage{listings}
\usepackage{xcolor}
\usepackage{booktabs}
\usepackage[margin=2.5cm]{geometry}
\usepackage{fancyhdr}

\hypersetup{
    colorlinks=true,
    linkcolor=blue,
    urlcolor=blue,
    citecolor=blue,
}

\definecolor{codebg}{rgb}{0.95,0.95,0.95}
\definecolor{codeframe}{rgb}{0.8,0.8,0.8}

\lstset{
    backgroundcolor=\color{codebg},
    frame=single,
    rulecolor=\color{codeframe},
    basicstyle=\ttfamily\small,
    breaklines=true,
    numbers=left,
    numberstyle=\tiny,
    keywordstyle=\color{blue},
    commentstyle=\color{green!40!black},
    stringstyle=\color{red},
    showstringspaces=false,
}

\pagestyle{fancy}
\fancyhf{}
\fancyhead[L]{ASTRA Plotter 使用手册}
\fancyhead[R]{\thepage}
\renewcommand{\headrulewidth}{0.4pt}

\title{\textbf{ASTRA Plotter 使用手册}\\[0.5em]
       \large 基于 Python 的 ASTRA 模拟数据后处理与可视化工具}
\author{ASTRA Plotter Team}
\date{\today}

\begin{document}

\maketitle
\tableofcontents
\newpage

% ============================================================================
\section{简介}
% ============================================================================

ASTRA Plotter 是对传统 BDplotter 代码的全面重构，提供模块化、可扩展的 ASTRA
（\textit{A Space charge TRacking Algorithm}）模拟数据后处理框架。

\subsection{功能特性}
\begin{itemize}
    \item \textbf{数据加载}：支持 ASTRA 输出的 Emit (.Xemit/.Yemit/.Zemit)、
          Sigma (.Sigma)、Cemit (.Cemit) 及相空间粒子文件
    \item \textbf{束流分析}：发射度计算、束斑统计、能散分析、电流分布、
          纵向切片分析、群聚因子
    \item \textbf{可视化}：出版级质量的发射度演化图、束斑尺寸图、能量曲线、
          相空间密度图（含投影）、切片参数综合图
    \item \textbf{交互式 GUI}：基于 ipywidgets 的交互式后处理界面
    \item \textbf{数据管理}：一键备份与时间戳归档
    \item \textbf{环境检测}：自动验证 Python 版本及依赖包兼容性
\end{itemize}

\subsection{与 BDplotter 的关系}
本项目继承了 BDplotter 的核心算法（发射度计算、本征发射度分析、相空间加载等），
并对代码进行了以下改进：
\begin{itemize}
    \item 模块化拆分（loader / analyzer / plotter / cosmetics）
    \item 规范化命名和文档字符串
    \item 全局配置集中管理（物理常数、文件格式 dtype）
    \item 增加 Jupyter Notebook 交互式界面
    \item 增加环境检测和备份工具
\end{itemize}

% ============================================================================
\section{安装与配置}
% ============================================================================

\subsection{系统要求}
\begin{itemize}
    \item 操作系统：macOS / Linux / Windows
    \item Python $\geq$ 3.7
    \item 可选：LaTeX 发行版（用于图表 TeX 渲染，如 MacTeX 或 TeX Live）
\end{itemize}

\subsection{安装步骤}

\subsubsection{方式 A：pip 安装}
\begin{lstlisting}[language=bash]
cd astra_plotter
pip install -r requirements.txt
\end{lstlisting}

\subsubsection{方式 B：Conda 环境}
\begin{lstlisting}[language=bash]
conda env create -f environment.yml
conda activate astra_plotter
\end{lstlisting}

\subsection{环境检测}
运行环境检测模块确认所有依赖就绪：
\begin{lstlisting}[language=bash]
python -m astra_plotter.utils.env_check
\end{lstlisting}

\begin{lstlisting}[language=python]
from astra_plotter.utils.env_check import check_environment
result = check_environment()
\end{lstlisting}

% ============================================================================
\section{项目结构}
% ============================================================================

\begin{lstlisting}[language=bash]
astra_plotter/
├── astra_plotter/          # 核心包
│   ├── __init__.py
│   ├── config.py           # 全局配置
│   ├── core/
│   │   ├── loader.py       # 数据加载
│   │   ├── analyzer.py     # 束流分析
│   │   ├── plotter.py      # 绘图
│   │   └── cosmetics.py    # 颜色映射与样式
│   └── utils/
│       ├── env_check.py    # 环境检测
│       └── helpers.py      # 工具函数
├── gui/                    # GUI 子项目
├── notebooks/              # Jupyter 笔记本
├── scripts/                # CLI 工具
├── data/                   # 备份存储
├── examples/               # 示例
├── tests/                  # 测试
└── docs/                   # 文档
\end{lstlisting}

% ============================================================================
\section{ASTRA 文件格式}
% ============================================================================

\subsection{Emit 文件 (.Xemit / .Yemit / .Zemit)}

记录了 RMS 束流参数沿纵向位置 $z$ 的演化。

\begin{table}[h]
\centering
\caption{Emit 文件列定义}
\begin{tabular}{@{}lll@{}}
\toprule
\textbf{列} & \textbf{名称} & \textbf{单位} \\
\midrule
1 & z      & m \\
2 & t      & ns \\
3 & $\langle x \rangle$ / $\langle p \rangle$ & m / eV/c \\
4 & $\sigma_x$ / $\sigma_t$ & m / ns \\
5 & $\sigma_{x'}$ / $\sigma_p$ & rad / eV/c \\
6 & $\gamma\epsilon_x$ / $\epsilon_z$ & m·rad / eV·s \\
7 & corr   & -- \\
\bottomrule
\end{tabular}
\end{table}

\subsection{Sigma 文件 (.Sigma)}

包含完整的 $6 \times 6$ 束团协方差矩阵（上三角元素），共 21 列：
\begin{equation}
\Sigma = \begin{pmatrix}
\langle x^2 \rangle & \langle x x' \rangle & \langle x y \rangle &
\langle x y' \rangle & \langle x z \rangle & \langle x z' \rangle \\
& \langle x'^2 \rangle & \langle x' y \rangle & \langle x' y' \rangle &
\langle x' z \rangle & \langle x' z' \rangle \\
& & \langle y^2 \rangle & \langle y y' \rangle & \langle y z \rangle &
\langle y z' \rangle \\
& & & \langle y'^2 \rangle & \langle y' z \rangle &
\langle y' z' \rangle \\
& & & & \langle z^2 \rangle & \langle z z' \rangle \\
& & & & & \langle z'^2 \rangle
\end{pmatrix}
\end{equation}

\subsection{相空间文件}

每行代表一个宏粒子，共 10 列：
$x$, $y$, $z$, $P_x$, $P_y$, $P_z$, clock, charge, index, status\_flag。

其中动量单位为 eV/c，坐标单位为 m。$z$ 和 $P_z$ 为累积值，
加载时自动转换为绝对值。仅保留 status\_flag $> 0$ 的存活粒子。

% ============================================================================
\section{API 参考}
% ============================================================================

\subsection{数据加载模块 (loader)}

\subsubsection{LoadAstraEmit(rootname, run)}
加载 Xemit、Yemit、Zemit（及可选的 Cemit）文件。

\textbf{参数：}
\begin{itemize}
    \item \texttt{rootname} (str): ASTRA 根文件名，如 \texttt{'astra'}
    \item \texttt{run} (str): 运行编号，如 \texttt{'001'}
\end{itemize}

\textbf{返回：} \texttt{(X, Y, Z, C)} 四个 NumPy 结构化数组。

\subsubsection{LoadAstraSigma(rootname, run)}
加载 Sigma 文件并计算本征发射度。

\textbf{返回：} \texttt{(S, enx, eny, enz)}，其中 enx/eny 为横向本征发射度，
enz 为纵向发射度。

\subsubsection{LoadAstraPhaseSpace(filename)}
加载相空间粒子文件，自动筛选存活粒子并转换坐标。

\textbf{返回：} 结构化数组，包含各粒子的相空间坐标和动量。

\subsection{束流分析模块 (analyzer)}

\subsubsection{Analysis(phsp)}
计算束团全局统计量。

\textbf{返回：} 字典，包含粒子数、平均 $\gamma$/$\beta$、
RMS 束斑尺寸、归一化发射度、相对能散等。

\subsubsection{UniformSliceAnalysis(phsp, numbins, bunch\_charge)}
等间隔纵向切片分析。

\textbf{参数：}
\begin{itemize}
    \item \texttt{phsp}: 相空间数组
    \item \texttt{numbins}: 切片数量
    \item \texttt{bunch\_charge}: 束团总电荷 (C)
\end{itemize}

\textbf{返回：} \texttt{(numbins, 14)} 矩阵，每行包含切片中心位置、
电流、发射度、亮度等参数。

\subsubsection{current(z, charge, Nbins)}
计算纵向电流分布 $I(z)$。

\subsubsection{BunchFormFactor(z, kmin, kmax, nk, IsLog)}
计算群聚因子 $|F(k)|^2$。

\subsection{绘图模块 (plotter)}

所有绘图函数接受 ASTRA 结构化数组，返回 \texttt{matplotlib.figure.Figure}。

\begin{table}[h]
\centering
\caption{绘图函数一览}
\begin{tabular}{@{}ll@{}}
\toprule
\textbf{函数} & \textbf{说明} \\
\midrule
\texttt{PlotEmit1plt(X, Y, Z)} & 三方向发射度演化（双 y 轴）\\
\texttt{PlotSize1plt(X, Y, Z)} & 束斑尺寸演化\\
\texttt{PlotEnergy1plt(X, Y, Z)} & 能量与能散演化\\
\texttt{PlotEigenEmits(S, enx, eny, enz)} & 本征发射度\\
\texttt{DensityPlot(X, Y, Nbins)} & 2D 密度图（对数色标）\\
\texttt{DensityplotwProjec2x2(X, Y, Nbins)} & 密度图 + 双投影\\
\texttt{DensityPlot\_w\_proj(X, Y, Nbins)} & 密度图 + 投影叠加\\
\texttt{PlotSliceParameters(sm)} & 切片分析综合图\\
\texttt{PlotBunchFormFactor(k, bff)} & 群聚因子图\\
\bottomrule
\end{tabular}
\end{table}

% ============================================================================
\section{使用示例}
% ============================================================================

\subsection{命令行快速绘图}
\begin{lstlisting}[language=bash]
# 切换到数据目录
cd simulation_output

# 绘制标准三图
python ../scripts/lineplot.py astra 001

# 仅绘制发射度，保存为 PDF
python ../scripts/lineplot.py astra 001 --type emit --save emit.pdf
\end{lstlisting}

\subsection{Jupyter Notebook 完整分析}
\begin{lstlisting}[language=python]
# 在 notebooks/astra_plotter.ipynb 中运行
from astra_plotter import *
from astra_plotter.core.cosmetics import set_publication_style

set_publication_style(font_size=14)

# 加载数据
X, Y, Z, C = LoadAstraEmit('astra', '001')
PhSp = LoadAstraPhaseSpace('astra.0500.001')

# 统计分析
result = Analysis(PhSp)
print(f"ε_nx = {result['nemit_x']*1e6:.2f} µm")

# 切片分析
slices = UniformSliceAnalysis(PhSp, 30, 1.0e-9)

# 绘制全套图表
PlotEmit1plt(X, Y, Z)
PlotSize1plt(X, Y, Z)
PlotEnergy1plt(X, Y, Z)
DensityplotwProjec2x2(PhSp['x']*1e3,
                      PhSp['px']/PhSp['pz'].mean(), 50)
plt.show()
\end{lstlisting}

\subsection{交互式 GUI}
\begin{lstlisting}[language=python]
from gui import PostProGUI, build_jupyter_gui

gui = PostProGUI(data_dir='simulation_output')
gui.load_data()
gui.show_statistics()

# 在 Jupyter 中显示交互式界面
widget = build_jupyter_gui(gui)
display(widget)
\end{lstlisting}

% ============================================================================
\section{物理常数与约定}
% ============================================================================

\begin{table}[h]
\centering
\caption{使用的物理常数}
\begin{tabular}{@{}lll@{}}
\toprule
\textbf{符号} & \textbf{数值} & \textbf{说明} \\
\midrule
$c$    & $2.99792458 \times 10^8$ m/s & 真空光速 \\
$m_e c^2$ & $0.5109989461 \times 10^6$ eV & 电子静止能量 \\
\bottomrule
\end{tabular}
\end{table}

\textbf{发射度约定：}
归一化发射度定义为
$\varepsilon_{n} = \beta\gamma \sqrt{\langle x^2\rangle\langle x'^2\rangle -
\langle xx'\rangle^2}$，
其中 $x' = P_x / P_z$ 为几何发散角。

\textbf{单位约定：}
\begin{itemize}
    \item 位置/束斑：输出时通常以 mm 表示
    \item 发射度：输出时通常以 µm 表示
    \item 动量：eV/c
    \item 时间：ns
\end{itemize}

% ============================================================================
\section{备份与数据管理}
% ============================================================================

备份脚本将 \texttt{simulation\_output/} 中的所有文件
复制到 \texttt{data/YYYYMMDD\_HHMMSS/} 目录。

\begin{lstlisting}[language=bash]
# 备份默认目录
python scripts/backup.py

# 备份指定目录
python scripts/backup.py /path/to/simulation

# 列出已有备份
python scripts/backup.py --list
\end{lstlisting}

备份目录中自动生成 \texttt{\_backup\_info.txt} 记录备份元信息。

% ============================================================================
\section{常见问题}
% ============================================================================

\subsection{LaTeX 渲染失败}
如果系统未安装 LaTeX，matplotlib 的 TeX 渲染将失败。
解决方法：
\begin{enumerate}
    \item macOS：安装 MacTeX (\texttt{brew install --cask mactex})
    \item Linux：安装 TeX Live (\texttt{sudo apt install texlive})
    \item 或者设置 \texttt{use\_tex=False} 禁用 TeX 渲染
\end{enumerate}

\subsection{找不到相空间文件}
确认 \texttt{simulation\_output/} 目录中存在格式为
\texttt{rootname.NNNN.NNN} 的文件。
可使用 \texttt{discover\_astra\_files()} 函数自动检测。

\subsection{切片分析结果异常}
切片分析需要足够的粒子统计量。建议每个切片至少包含 10 个粒子。
可以使用 \texttt{UniformSliceAnalysis} 的 \texttt{numbins} 参数调节切片数量。

% ============================================================================
\section{开发指南}
% ============================================================================

\subsection{添加新数据格式}
\begin{enumerate}
    \item 在 \texttt{astra\_plotter/config.py} 中定义 dtype
    \item 在 \texttt{astra\_plotter/core/loader.py} 中添加加载函数
    \item 在 \texttt{astra\_plotter/\_\_init\_\_.py} 中导出
\end{enumerate}

\subsection{添加新绘图类型}
\begin{enumerate}
    \item 在 \texttt{astra\_plotter/core/plotter.py} 中添加函数
    \item 函数应返回 \texttt{matplotlib.figure.Figure}
    \item 在 \texttt{gui/\_\_init\_\_.py} 的 \texttt{PostProGUI.plot()} 中注册
\end{enumerate}

% ============================================================================

\appendix
\section{SLAC-DESY 颜色映射}

本项目使用 SLAC-DESY 标准束流颜色映射（由 C. Behrens, DESY 和
T. Maxwell, SLAC 开发）。该颜色映射专为束流动力学可视化优化，
在低密度区域呈现深色、高密度区域呈现暖色，具有良好的感知均匀性。

\end{document}
